\section{Related Work}
\label{sec:related-work}

\paragraph{Hyperbolic representations and learned curvature.}
Poincar\'{e} embeddings use negative curvature for compact hierarchical representations, while Lorentz embeddings provide an alternative model with closed-form geodesic operations~\citep{nickel2017poincare,nickel2018lorentz}. Learning curvature is already established: product-space embeddings jointly learn geometry across factors~\citep{gu2019mixed}, and HGCN learns a curvature at each graph-convolution layer~\citep{chami2019hgcn}. For knowledge graphs, MuRP provides a hyperbolic relational model~\citep{balazevic2019poincare}, while RotH and AttH learn relation-specific curvatures together with isometric transformations~\citep{chami2020atth}. Poincar\'{e} GloVe extends hyperbolic representation learning to lexical co-occurrence~\citep{tifrea2019poincare}. Grouped or relation-specific curvature is therefore a precedent, rather than a novelty claim of this paper.

\paragraph{Riemannian optimization and optimizer memory.}
RSGD establishes stochastic optimization directly on a manifold~\citep{bonnabel2013sgd}; curvature-dependent first-order analysis characterizes geodesically convex objectives~\citep{zhang2016geodesic}. Riemannian Adam and AMSGrad maintain adaptive statistics on a product manifold and transport first-moment information between tangent spaces~\citep{becigneul2019radam}. More recent adaptive frameworks analyze minibatch size and constant or diminishing steps on embedded manifolds~\citep{sakai2024adaptive}. These are important fixed-manifold foundations. Our analysis additionally tracks a learned curvature coordinate, including the gradient induced by identifying points through a curvature-dependent homothety. It retains the different representation and curvature step weights instead of assigning both blocks the same rate by default.

\paragraph{Curvature-aware state mapping: the closest optimization precedent.}
\citet{bdeir2025robust} already identify mismatches caused by updating curvature while retaining stale manifold parameters and optimizer states. Their curvature-aware scheme updates parameters and transports gradients and momentum through the origin tangent space; it is combined with Riemannian AdamW and radius rescaling. We therefore do not claim to introduce cross-curvature optimizer-state consistency. HyperbolicMomentum instead studies a homothetic point-and-momentum map, grouped slow curvature updates, and a curvature--momentum normalized geodesic step. The distinction is the chosen coupled update and its error-explicit analysis; the homothety itself also has precedents discussed by that work. Its behavior relative to origin-tangent remapping requires a matched empirical comparison, which the current experiments do not supply.

\begin{table*}[!ht]
\centering
\caption{\textbf{Mechanisms in related work.} Keywords identify the design; checks identify properties of the cited formulation. Curvature learning and cross-curvature optimizer-state mapping both have prior art. The final column denotes our specific curvature--momentum step normalization, not all forms of adaptive learning rates or radius control.}
\label{tab:related-work-features}
\begingroup
\footnotesize
\setlength{\tabcolsep}{4.2pt}
\renewcommand{\arraystretch}{1.12}
\begin{tabularx}{\linewidth}{@{}>{\raggedright\arraybackslash}p{.215\linewidth}>{\raggedright\arraybackslash}p{.225\linewidth}*{5}{>{\centering\arraybackslash}X}@{}}
\toprule
\textbf{Method / family} & \textbf{Core mechanism} & \textbf{Learned $K$} & \textbf{$K$ scope} & \textbf{Tangent memory} & \textbf{Cross-$K$ memory} & \textbf{Step control} \\
\midrule
\rowcolor{rwgray}
\multicolumn{7}{@{}l}{\textit{Optimization on a prescribed geometry}} \\
RSGD~\citep{bonnabel2013sgd} & Stochastic tangent descent & \rwcross & Fixed & \rwcross & \rwna & \rwcross \\
Lorentz embeddings~\citep{nickel2018lorentz} & Hyperboloid Exp update & \rwcross & Fixed & \rwcross & \rwna & \rwcross \\
RAdam / RAMSGrad~\citep{becigneul2019radam} & Transported adaptive moments & \rwcross & Fixed & \rwcheck & \rwna & \rwcross \\
\addlinespace[2pt]
\rowcolor{rwgray}
\multicolumn{7}{@{}l}{\textit{Learning the representation geometry}} \\
Product spaces~\citep{gu2019mixed} & Mixed-curvature factors & \rwcheck & Factor & \rwna & \rwna & \rwcross \\
HGCN~\citep{chami2019hgcn} & Hyperbolic message passing & \rwcheck & Layer & \rwna & \rwna & \rwcross \\
RotH / AttH~\citep{chami2020atth} & Isometries + attention & \rwcheck & Relation & \rwcross & \rwcross & \rwcross \\
\addlinespace[2pt]
\rowcolor{rwgray}
\multicolumn{7}{@{}l}{\textit{Coupling curvature changes to optimizer state}} \\
Curvature-aware Lorentz~\citep{bdeir2025robust} & Origin-tangent remapping + RAdamW & \rwcheck & Manifold & \rwcheck & \rwcheck & \rwcross \\
\midrule
\rowcolor{oursbg}
\textbf{HyperbolicMomentum} & \textbf{Homothety + controlled Exp step} & \rwcheck & Group & \rwcheck & \rwcheck & \rwcheck \\
\bottomrule
\end{tabularx}
\par\vspace{4pt}
\begin{minipage}{\linewidth}
\scriptsize
\rwcheck\ present; \rwcross\ absent in the listed formulation; \rwna\ not specified as a core optimizer mechanism, or not applicable for fixed $K$.
\emph{Tangent memory}: a first-moment state attached to moving manifold points; ordinary Euclidean Adam states in origin coordinates are not counted.
\emph{Cross-$K$ memory}: an explicit rule for adapting that state when curvature changes; mappings of neural activations between layers are not counted.
\emph{Normalized step}: $\eta^{\rm geo}=\eta/(1+\gamma\sqrt K\|\widetilde m\|)$, which bounds $\sqrt K\|v\|$; Adam preconditioning and embedding-radius clipping are distinct mechanisms.
The row for our method describes its grouped-manifold formulation, not the historical shared-entity KG implementation. No check implies empirical superiority.
\end{minipage}
\endgroup
\end{table*}


\paragraph{The map being compared.}
The origin-tangent point map has the form
$\Exp_{o_{K'}}^{K'}\!\circ\Log_{o_K}^{K}$ and retains the physical radial coordinate. Our map is
$\Psi_{K\to K'}(x)=\sqrt{K/K'}\,x$, with differential
$D\Psi(v)=\sqrt{K/K'}\,v$; it retains the canonical state
$q=\sqrt Kx$, $p=\sqrt Km$ while scaling physical distances.
These invariants are different. Neither map implies a universal performance ordering. Table~\ref{tab:related-work-features} separates architectural curvature learning, transported tangent memory, and changes of the geometry itself, rather than assigning every hyperbolic model the same optimizer.

\paragraph{Generalization and evidence boundaries.}
Replace-one stability links coupled optimization trajectories to expected generalization under an appropriate sampling model~\citep{hardt2016stability}. Our conditional bound exposes its curvature-radius and cumulative-step dependence and distinguishes fixed i.i.d. data from graph dependence or continual-learning shift. Historical KG measurements in this paper are exploratory fixed-versus-learned-curvature Adam comparisons; they are not matched optimizer comparisons against RotH/AttH or curvature-aware RAdamW. The table records mechanisms, while the experiment section separately records which comparisons were actually run.
